% Generated draft: requires mathematical and source review.
\documentclass[11pt]{article}
\usepackage[margin=25mm]{geometry}
\usepackage{fontspec}
\setmainfont{Latin Modern Roman}
\usepackage{amsmath,amssymb,mathtools,unicode-math}
\setmathfont{Latin Modern Math}
\usepackage{enumitem,needspace,xurl,hyperref}
\setlength{\parindent}{0pt}
\setlength{\parskip}{5pt}
\setlength{\emergencystretch}{3em}
\newcommand{\N}{\mathbb{N}}
\newcommand{\Z}{\mathbb{Z}}
\newcommand{\Q}{\mathbb{Q}}
\newcommand{\R}{\mathbb{R}}
\newcommand{\C}{\mathbb{C}}
\newcommand{\F}{\mathbb{F}}
\newcommand{\A}{\mathbb{A}}
\newcommand{\PP}{\mathbb{P}}
\newcommand{\E}{\mathbb{E}}
\newcommand{\eps}{\varepsilon}
\newcommand{\dd}{\,\mathrm{d}}
\begin{document}
\section*{UnsolvedMath ID: 20001949}
\textbf{Source problem code:} catalogue entry
\subsection{Definitions and mathematical statement}
Let \((X,g,I,J,K)\) be a hyperk\"ahler manifold with a hyperHamiltonian action of a compact Lie group \(G\). Choose an invariant inner product on \(\mathfrak g\), and let
\[
 \mu_{\mathrm{HK}}=(\mu_I,\mu_J,\mu_K):X\to\mathfrak g^*\otimes\R^3,
 \qquad f=\|\mu_{\mathrm{HK}}\|^2.
\]
Suppose that for every \(x\in X\), its forward negative-gradient trajectory for \(f\) is defined for all \(t\geq0\) and lies in a compact subset of \(X\), allowed to depend on \(x\). Put \(Z=\mu_{\mathrm{HK}}^{-1}(0)\). The reconstructed AIM conjecture asks whether inclusion induces a surjection
\[
 H_G^*(X;R)\longrightarrow H_G^*(Z;R),
 \qquad H_G^*(Y;R)=H^*(EG\times_GY;R).
\]
The original calls the acting group \(G\) but indexes cohomology by an undefined \(K\); the aid explicitly sets that \(K=G\). It also supplies no coefficient ring \(R\), so coefficient-dependent versions must be distinguished.

Comment \(1.2\) lists compact \(X\), proper \(f\), and the above trajectory-compactness condition, stating the implications compact \(X\Rightarrow\) proper \(f\Rightarrow\) compact forward trajectories, and observing that the third is the relevant flexible hypothesis in the intended applications. Its attribution to Kirwan and claim that the first two almost never hold are source commentary; they are not additional independent hypotheses, and the comment ends mid-sentence.

The aid reports that this exact flow-closed implication is known over \(\Q\) for compact tori, while the nonabelian version remains unresolved in literature checked through 9 August 2026. It distinguishes both quiver-variety quotient-surjectivity theorems and nonabelian counterexamples in other settings from this precise level-set statement; no unqualified inference between them is made.

Its conditional coefficient-sensitive criterion is the following. Suppose the full function \(f\) provides a countable equivariant Morse--Kirwan exhaustion
\[
 X_0\subset X_1\subset\cdots,\qquad \bigcup_jX_j=X,
\]
whose bottom \(X_0\) equivariantly retracts to \(Z\). Assume each attachment has an equivariant Thom description by an \(R\)-oriented negative bundle \(\nu_j^-\) of rank \(\lambda_j\) over its critical set \(C_j\), with the relevant Thom map restricting to multiplication by the Euler class
\(e_G(\nu_j^-)\). Require that this multiplication be injective on \(H_G^*(C_j;R)\). Then the equivariant perfection step yields degreewise short exact sequences
\[
 0\longrightarrow H_G^{*-\lambda_j}(C_j;R)
 \longrightarrow H_G^*(X_j;R)
 \longrightarrow H_G^*(X_{j-1};R)\longrightarrow0.
\]
Finally assume the exhaustion/Borel construction satisfies the continuity conditions for the Milnor sequence
\[
 0\longrightarrow\varprojlim{}^{1}H_G^{k-1}(X_j;R)
 \longrightarrow H_G^k(X;R)
 \longrightarrow\varprojlim H_G^k(X_j;R)\longrightarrow0.
\]
The surjective transition maps make the derived-limit term vanish, and the inverse-limit projection onto the bottom is surjective. Under all these hypotheses the desired restriction to \(Z\) is surjective. Countability, negative-bundle orientation, Euler non-zero-divisibility, bottom retraction and continuity at infinity are substantive assumptions, not consequences established here from trajectory compactness alone.

For arbitrary invariant functions that weaker dynamical condition is insufficient: on \(S^2\), with trivial group and \(f=z^2\), trajectories stay compact while \(f^{-1}(0)\) is the equator \(S^1\), and \(H^1(S^2;\Q)=0\) cannot surject onto \(H^1(S^1;\Q)\). This is the aid's diagnostic for missing Euler/perfection input; it is not a hyperk\"ahler moment-map counterexample.
\subsection{Short English statement}
The reconstructed conjecture asks whether compact forward trajectories of the full hyperk\"ahler moment-map norm-square imply equivariant surjectivity to its zero level. The aid reports the rational torus case and supplies a conditional Morse-exhaustion criterion, including Euler and inverse-limit requirements. It does not prove those requirements in the general nonabelian case.
\subsection{Sources}
\begin{itemize}
\item Local UnsolvedMath registry, permanent ID \(20001949\), source code \texttt{catalogue entry}; canonical collection \texttt{aim-geometry-notes.json}, notes[286]. Statement status \texttt{reconstructed\_unverified}; clean statement null.
\item AIM workshop ``Moment maps and surjectivity in various geometries,'' Conjecture \(1.1\) and truncated Comment \(1.2\): \url{https://aimath.org/WWN/momentmaps/momentmaps.pdf}. The undefined cohomology-index group \(K\) is explicitly reconstructed as the acting \(G\). No PDF was opened in this drafting session.
\item Kirwan is named in the comment. No full authoritative citations for the torus theorem, quiver-variety comparison results, nonabelian examples or Morse--Kirwan inputs are supplied; original remarks and literature fields are empty, with no proposer or year.
\item Supplied research attempt \(1\), result ``reduction,'' run status \texttt{partially\_solved}, rigor label ``proved,'' novelty confidence ``low.'' Registry classification \texttt{PARTIAL-PROGRESS}, suggested difficulty \texttt{L4}, updated 14 August 2026. The nonabelian literature qualification is checked only through 9 August 2026.
\end{itemize}
\end{document}
